\section{Analiza i przygotowanie danych}
\label{sec:dane}

\subsection{Statystyki zbioru}

Surowy zbiór liczy \textbf{44\,898} artykułów o niemal zrównoważonych klasach
($\sim$52\% Fake / 48\% Real). Eksploracyjna analiza danych
(\texttt{notebooks/explore/01\_eda.ipynb}) ujawniła dwa istotne problemy jakości:

\begin{itemize}[noitemsep]
  \item \textbf{Duplikaty}: ok. 13{,}9\% artykułów to dokładne duplikaty tekstu,
        skupione niemal wyłącznie w klasie Fake (ok. 25{,}7\% artykułów Fake vs
        1{,}1\% Real). Brak konfliktów etykiet między klasami --- deduplikacja jest bezpieczna.
  \item \textbf{Puste / bardzo krótkie teksty}: ok. 1{,}4\% rekordów (głównie wpisy
        Fake zawierające jedynie odnośnik do YouTube).
\end{itemize}

Po pełnym potoku czyszczenia i deduplikacji pozostaje \textbf{38\,470} artykułów
(Fake: 17\,281 = 44{,}9\%, Real: 21\,189 = 55{,}1\%). Rozkład długości tekstu jest
silnie prawoskośny; 95.~percentyl to ok. \textbf{900 słów} --- wartość przyjęta jako
kandydat na \texttt{max\_seq\_len} w docelowym modelu sekwencyjnym.

\subsection{Wyciek informacji (data leakage)}

Analiza w \texttt{notebooks/explore/02\_leakage.ipynb} wykorzystała informację
wzajemną (Mutual Information, MI) oraz test $\chi^2$ do oceny siły poszczególnych
markerów jako predyktorów klasy. Najsilniejsze sygnały (Tabela~\ref{tab:leakage})
to artefakty redakcyjno-techniczne, a nie cechy treści: nagłówek agencyjny
``\texttt{(Reuters)}'' pozwala niemal jednoznacznie wskazać klasę Real.

\begin{table}[H]
  \centering
  \caption{Najsilniejsze markery wycieku wg informacji wzajemnej (MI).}
  \label{tab:leakage}
  \begin{tabular}{lcl}
    \toprule
    \textbf{Marker} & \textbf{MI [nat]} & \textbf{Dominująca klasa} \\
    \midrule
    słowo ,,reuters''            & 0{,}650 & Real \\
    dateline ,,MIASTO (Reuters) -- '' & 0{,}456 & Real \\
    ,,featured image''           & 0{,}136 & Fake \\
    \texttt{@handle} (media społ.)& 0{,}076 & Fake \\
    ,,getty (images)''           & 0{,}061 & Fake \\
    prefiks clickbait (VIDEO/WATCH) & 0{,}018 & Fake \\
    \bottomrule
  \end{tabular}
\end{table}

Wniosek z analizy: kolumna \texttt{subject} rozdziela klasy w 100\% (każda kategoria
należy w całości do jednej klasy), a \texttt{title} oraz \texttt{date} również niosą
silny wyciek (Fake: 83\% tytułów ALL-CAPS; daty Fake w niestandardowym formacie).
Wszystkie trzy kolumny są więc odrzucane --- model korzysta wyłącznie z \texttt{text}.

\subsection{Czyszczenie tekstu}

Funkcja \texttt{clean\_text} (\texttt{src/data/cleaning.py}) stosuje uporządkowaną
kaskadę wyrażeń regularnych --- każdy wzorzec ma uzasadnienie w analizie MI powyżej.
Kolejność operacji jest istotna (np.\ usunięcie ,,via @handle'' przed samym
\texttt{@handle}, rozdzielenie ,,2017World'' przed sprowadzeniem do małych liter):

\begin{enumerate}[noitemsep]
  \item markery agencyjne Reuters (dateline, ,,(Reuters)'', słowo ,,reuters'');
  \item adresy URL (\texttt{http(s)}, \texttt{www}, \texttt{pic.twitter.com}, \texttt{tmsnrt.rs});
  \item uchwyty mediów społecznościowych (\texttt{@handle}, ,,via @\dots'');
  \item prefiksy clickbait (\texttt{VIDEO:}, \texttt{WATCH:}, \texttt{BREAKING:});
  \item szablony podpisów zdjęć (,,featured image via getty'', ,,photo by\dots'');
  \item artefakty HTML (znaczniki, encje, bloki \texttt{<script>});
  \item sygnatury serwisów (,,21st Century Wire'');
  \item normalizacja: białe znaki, usunięcie pojedynczych znaków, małe litery.
\end{enumerate}

Po czyszczeniu stosowany jest filtr \texttt{filter\_short\_articles} z progiem
\textbf{\texttt{min\_words = 10}} (usuwa głównie wpisy-zaślepki Fake). Walidacja w
\texttt{04\_post\_cleaning.ipynb} potwierdza, że kluczowe markery (dateline, URL,
\texttt{@handle}, getty) spadają do 0\% wystąpień, a pokrycie słownika GloVe pozostaje
wysokie (OOV $\le$ 0{,}2\%).

\subsection{Podział na zbiory}

Funkcja \texttt{make\_splits} (\texttt{src/data/splits.py}) dzieli dane
\textbf{stratyfikowanie} po etykiecie na zbiory train/val/test w proporcji
60/20/20 (dwukrotne wywołanie \texttt{train\_test\_split}), ze stałym ziarnem
\texttt{RANDOM\_STATE} dla powtarzalności (Tabela~\ref{tab:splits}). Balans klas
($\sim$44{,}9\% Fake / 55{,}1\% Real) jest zachowany we wszystkich trzech zbiorach.

\begin{table}[H]
  \centering
  \caption{Podział stratyfikowany na zbiory.}
  \label{tab:splits}
  \begin{tabular}{lcc}
    \toprule
    \textbf{Zbiór} & \textbf{Liczba artykułów} & \textbf{Udział} \\
    \midrule
    train & 23\,082 & 60\% \\
    val   & 7\,694  & 20\% \\
    test  & 7\,694  & 20\% \\
    \midrule
    razem & 38\,470 & 100\% \\
    \bottomrule
  \end{tabular}
\end{table}
